
% ====================================================================
% The pinned-component proposition: deriving the attractor size from
% structure alone. Data check: experiments/attractor_theorem.py
% ====================================================================
\subsection{A derivation: the pinned component from structure alone}
\label{sec:attractor-thm}

The previous subsection measured the pinned component and showed it tracks
tube/sp. We now derive its size from the graph alone, with no simulation, and
check the derivation against the measurement.

\paragraph{Setup.} Fix a graph $G$ and demand $D$. A router maps each demand
pair $(i,j)$ to a unit flow along $i$--$j$ paths, placing load $\ell_r(e)$ on
edge $e$. Two routers agree on $e$ to the extent that the flow there is forced
rather than chosen. We make ``forced'' precise per pair.

\paragraph{Per-pair pinning.} Take one pair $(i,j)$ with shortest-path length
$d_{ij}$, and let $T_{ij}$ be the edges lying on some path of length at most
$\alpha\,d_{ij}$---the $\alpha$-stretch tube the tube/sp ratio counts. Any
reasonable router confines the $(i,j)$ flow to $T_{ij}$. If $|T_{ij}|=d_{ij}$
the tube is a single path: the flow has nowhere to go but those $d_{ij}$ edges
and every router, blind or not, loads them identically. If $|T_{ij}|>d_{ij}$
the pair has $|T_{ij}|/d_{ij}$ alternatives over which routers may disagree.
The forced fraction of the pair's flow is therefore
$\pi_{ij}=d_{ij}/|T_{ij}|$: one for a bare path, shrinking as the tube widens.

\paragraph{Aggregate.} The pinned component of the whole load distribution is
the flow-weighted mean of the per-pair pinning, weighting each pair by the
mass it moves (proportional to $d_{ij}$ under uniform demand):
\begin{equation}
\widehat{\pi}(G) \;=\;
\frac{\sum_{i<j} d_{ij}\,\pi_{ij}}{\sum_{i<j} d_{ij}}
\;=\;
\frac{\sum_{i<j} d_{ij}^{\,2}/|T_{ij}|}{\sum_{i<j} d_{ij}}.
\label{eq:pinned}
\end{equation}
This is a pure graph functional: no capacities, no traffic, no router. It is
also, by construction, a decreasing function of tube/sp---wider tubes mean
smaller $\pi_{ij}$---so the same structural axis that predicts routing gain
predicts the attractor's strength.

\paragraph{The derivation matches the measurement.} If the argument is right,
$\widehat{\pi}(G)$---computed from the graph---should predict the
physics-versus-blind similarity we actually measured by simulation. It does:
across the fifteen topologies, $\widehat{\pi}$ correlates with the measured
phys--ECMP cosine at Spearman $+0.86$ (Pearson $+0.80$, $p<10^{-3}$), a
structure-only quantity recovering a simulated one. The relationship is
concave, not linear: as $\widehat{\pi}\!\to\!1$ the cosine saturates against its
ceiling of one, since a similarity bounded above cannot grow without limit. The
fit leaves a floor near $0.77$ at $\widehat{\pi}=0$---the residual agreement two
routers retain from shared graph geometry even when nothing is pinned.

\paragraph{What this is, and is not.} Equation~\eqref{eq:pinned} turns the
attractor from an observed coincidence into a quantity derivable from the graph:
the load distribution has a forced part, fixed by the tube geometry of the
demand pairs, and its size is $\widehat{\pi}(G)$. This is a derivation of the
attractor's \emph{strength}, not yet a closed bound on the per-router error:
$\pi_{ij}$ counts tube edges rather than solving the min-cut that would bound
exactly how much two routers can differ, and the concave link from
$\widehat{\pi}$ to cosine is fitted, not derived in closed form. What we claim is
narrower and, we think, sufficient: the attractor's size is a graph functional,
it is computable without simulating anything, and it predicts the measured
convergence. The four-way coincidence of Section~\ref{sec:dissociation} has a
single root we can now write down.
